\documentclass{article}
\usepackage{/globals/boxes}
\usepackage{/globals/tables}
\usepackage{/globals/anki}
\ankiprefix{schaltwerke}
\ankitopic{Schaltwerke}
\ankishow
\begin{document}
\section{Schaltwerke} 
Ein Schaltwerk ist ein Schaltnetz mit eine Rückkopplung; ein Ausgang ist gleichzeitig auch ein Eingang.
\ankinote{def}{Definition}{Schaltnetz mit Rückkopplung}
\begin{description}
 \item[Asynchron] Der Zustand des Schaltwerks kann sich konstant und unabhängig von allem anderen verändern.
 \item[Synchron] Der Zustand aller Schaltwerke verändern sich synchron, zentral gesteuert durch einen Takt. 
 \begin{description}
  \item[Pegelsteuerung] Der Zustand kann sich verändern wenn das Taktsignal einen bestimmten Wert hat. 
  \item[Flankensteuerung] Der Zustand kann sich das Taktsignal (zu einem bestimmten Wert) verändert.
 \end{description}
\end{description}
\ankinote{arten}{Arten}{Asynchron (immer veränderbar), Synchron (Takt)}
\ankinote{sync-arten}{Synchron Arten}{Pegelgesteuert (Wert), Flankensteuerung (Veränderung)}

\subsection{Flip-flops}
Ein Flip-flops kann jeweils ein Bit an Information speichern.
\begin{description}
\item[RS-flip-flop] Ein RS-flip-flop hat zwei Eingänge; dem Namen nach $R$ (Reset) und $S$ (Set), welche auch genau dies mit dem gespeicherten Wert tun.
\begin{ztable}{cc|c}{$R$ & $S$ & $Q$}
 0 & 0 & unverändert \\
 0 & 1 & 1 \\
 1 & 0 & 0 \\
 1 & 1 & ungültig \\
\end{ztable}
\ztodo{tikz}
\ankinote{rs-ff}{RS-flip-flop}{Reset, Set den Wert}

\item[Pegelgesteuerter D-flip-flop] Ein pegelgesteuerter D-flip-flop hat zwei eingänge; das Signal $D$ und und den Takt (die clock) $C$. Wenn $C=1$ wird $D$ gespeichert.
\begin{ztable}{cc|c}{$C$ & $D$ & $Q$}
 0 & 0 & unverändert \\
 0 & 1 & unverändert \\
 1 & 0 & 0 \\
 1 & 1 & 1 \\
\end{ztable}
\ztodo{tikz}
\ankinote{peg-d-ff}{Pegelgesteuerter D-flip-flop}{Wenn C=1, Wert=D}

\item[Flankengesteuerter D-flip-flop] Ein flankengesteuerter D-flip-flop hat die gleichen Eingänge wie ein pegelgesteuerter D-flip-flop, speichert aber den Wert von $D$ wenn $C$ von 1 zu 0 springt.
\begin{ztable}{cc|c}{$C$ & $D$ & $Q$}
 $1 \rightarrow{} 0$ & 0 & unverändert \\
 $1 \rightarrow{} 0$ & 1 & unverändert \\
 $0 \rightarrow{} 1$ & 0 & 0 \\
 $0 \rightarrow{} 1$ & 1 & 1 \\
\end{ztable}
\ztodo{tikz}
\ankinote{flank-d-ff}{Flankengesteuerter D-flip-flop}{Wenn C 0->1, Wert=D}
\end{description}
\end{document}